\subsection{Jeres band}
\subsubsection{Elsk din besætning}
Som revyband er der ingen grænser for hvilke genrer musik man bliver bedt om at spille.
Det kan være en udfordring i forhold til instrumentering.
Skulle man gengive hver genre autentisk ville det kræve flere instrumenter end der kan være på bandscenen, og flere musikere end en studenterrevy kan mønstre.
Det er selvfølgelig et fint udgangspunkt at forsøge at genskabe lyden af originalen, men nogle gange er det et dødfødt projekt.
For eksempel lader de fleste former for "computermusik" sig ikke oversætte direkte til instrumentering.
Hvis man insisterer alligevel, bliver man skuffet: musikere kan aldrig lave computermusik bedre end en computer.
Derfor er det vigtigt, når man bliver bedt om at spille et nummer, der ikke umiddelbart er skrevet til den besætning man har i bandet, at kunne \textit{fortolke} nummeret til ens besætning.
Med at fortolke nummeret menes at identificere nummerets essentielle elementer og bibeholde disse i ens arrangemnet, imens man tilpasser resten, så det fungerer med et band.
Målet er, at publikum nemt kan genkende nummeret, samtidig med at det lyder naturligt for bandet at spille det.

\subsubsection{Udvikl jeres stil}
Vær ikke bange for at genbruge idéer fra tidligere arrangementer. Over tid samler I erfaring med hvad der fungerer for jer I forskellige genrer. Læn jer ind i de greb I finder og gør det til jeres stil.

\subsubsection{Kend dit band}
Revybands samler ofte musikere med meget forskellig musikalsk baggrund. Det er  en styrke når man skal spille mange forskellige genrer.
Lær hinanden at kende musikalsk og udnyt hinandens styrker når I skriver arrangementer.
I stedet for at arrangere for instrumenter så prøv at tænke, at I arrangerer for bandet og musikerne i det.